Since the first direct detection of gravitational waves from a binary black hole merger~\citep{abbott2016gw150914}, deep-learning approaches to gravitational-wave parameter estimation have matured rapidly, from early proof-of-concept classifiers~\citep{george2018deep} to amortized neural posterior estimators that approach the accuracy of stochastic-sampling pipelines~\citep{gabbard2022vitamin,dax2021dingo} at a fraction of their latency.
Most of this progress has concentrated on parameters that imprint strongly on the observed strain, i.e., the chirp mass~\citep{cutler1994gw}, the merger time and signal-to-noise ratio~\citep{allen2012findchirp}, and the sky position~\citep{fairhurst2009triangulation,singer2016bayestar}.

Two extrinsic parameters have received far less attention as direct regression targets: the orbital phase at coalescence, $\varphi_c$, and the polarization angle, $\psi$.
There is a physical reason for this neglect.
For the dominant quadrupole ($\ell=2$, $|m|=2$) mode of a quasi-circular compact binary, the two polarization states are~\citep{sathyaprakash2009physics}
\begin{equation*}
    \begin{aligned}
    h_+ &\propto (1+\cos^2\iota)\cos\left(2\varphi(t)+2\varphi_c\right),\\
    h_\times &\propto 2\cos\iota\sin\left(2\varphi(t)+2\varphi_c\right),
    \end{aligned}
\end{equation*}
and the detector response mixes them through antenna patterns that depend on $2\psi$~\citep{cutler1994gw,sathyaprakash2009physics}.
Both $\varphi_c$ and $\psi$ enter the strain doubled, so the analytically degenerate combination in the face-on limit ($\cos\iota \to \pm 1$) is $2\varphi_c \pm 2\psi$, with the sign set by the sign of $\cos\iota$.
%The correction traces to how \texttt{coa\_phase} is injected for this population: it is the orbital reference phase in the convention used by the generator, and for a dominant-mode-only waveform that orbital phase enters the strain through $e^{i\cdot 2\cdot\text{coa\_phase}}$, exactly like $\psi$ already does.
In the face-on limit the signal becomes circularly polarized and depends on $\varphi_c$ and $\psi$ only through this single, corrected combination, so the individual angles are exactly unidentifiable there, and only partially disentangled as the inclination moves toward edge-on.
Bayesian samplers absorb this structure into broad, multimodal joint posteriors~\citep{veitch2010bayesian,aasi2013pe,veitch2015lalinference,ashton2019bilby}; a point-estimate regression network has no such refuge.

This paper asks the question:
\begin{quote}
    \itshape Can $\varphi_c$ and $\psi$ be recovered from strain alone by a deep neural network, individually, through their sum/difference combinations, or is the degeneracy effectively exact for a realistic detection-scale population?
\end{quote}

The answer, is a \textbf{null result}: across the five trunk architectures, two head parameterizations, and three values of a magnitude-regularization coefficient.
No model could perform measurably better than random guessing on either angle, while the same networks simultaneously recovered chirp mass, merger time, signal-to-noise ratio, and in some cases, sky position with high fidelity from the same shared features.

A null result of this kind is only as strong as the effort spent trying to break it.
The scientific contribution of this paper is therefore threefold.
First, the null itself, tested against an explicit list of confounds: data-pipeline faults, loss-wiring faults, architecture-specific artifacts, statistical flukes, and population heterogeneity.
Second, a methodological account of \emph{how} the null was established when the standard interpretability gate never cleared --- because unlike a gate failing for one model out of several, here it failed for every model at every regularization strength tried, forcing an explicit argument for why a broader, ungated battery of evidence should be trusted in its place, a response, at engineering-loop scale, to the same replication crisis that motivates pre-registration in quantitative science more broadly~\citep{ioannidis2005false,nosek2018preregistration}.
Third, a methodological lesson demonstrated directly within this study: a checkpoint that cleared every numeric threshold of the pre-registered interpretability gate turned out, on scatter-plot verification, to be a literal two-point mode collapse --- a passing aggregate metric that is not, by itself, proof of the property it is meant to certify.

%This null result carries a direct architectural implication, stated here rather than only in the Discussion.
%Since strain alone does not recover $\varphi_c/\psi$ under point estimation, but \S\ref{sec:phicpsi:prereq}'s analytic study shows the well-constrained combination \emph{is} recoverable given inclination, the natural next step is to condition the network on $\iota$ directly --- scoped as future work in \S\ref{sec:phicpsi:future}.

The research paper is organized as follows;
Section~\ref{sec:phicpsi:problem} formalizes the problem and the circular-regression framework.
Section~\ref{sec:phicpsi:prereq} presents the analytic prerequisite study, rerun under the corrected antenna-pattern model.
Section~\ref{sec:phicpsi:methods} details the datasets, architectures, losses, and training protocol, including the corrected combination construction.
Section~\ref{sec:phicpsi:arc} summarizes, briefly, two optimization pathologies in this training pipeline and how they were fixed prior to the main sweep.
Section~\ref{sec:phicpsi:results} presents the defended null result, its verification battery, and findings.
Section~\ref{sec:phicpsi:prereg} describes the pre-registered magnitude-penalty retune and its outcome.
Section~\ref{sec:phicpsi:discussion} discusses interpretation, limitations, and future work; Section~\ref{sec:phicpsi:conclusion} concludes.
